\chapter{Discrete-Time Market Model}
\label{chap:discrete_time_market_model}

In to financial modeling, transitioning from a single-period model to a multi-period model is essential for capturing the dynamic nature of trading. In this chapter, we lay out the mathematical foundations of a discrete-time multi-period market model.

\section{The Multi-Period Model Setup}

We consider a financial market consisting of one riskless asset (often called the money market account or a bank account) and $d \in \mathbb{N}$ risky assets (such as stocks or indices). Trading occurs at discrete time steps $t \in \{0, 1, \dots, T\}$, where $T \in \mathbb{N}$ denotes the finite time horizon of the model.

To formalize the uncertainty inherent in financial markets, we introduce a probability space $(\Omega, \mathcal{F}, \mathbb{P})$. 
\begin{itemize}
    \item $\Omega$ is the sample space, representing the set of all possible market scenarios (or paths of asset prices) up to time $T$.
    \item $\mathcal{F}$ is the $\sigma$-algebra of events, representing all possible observable sets of scenarios.
    \item $\mathbb{P}$ is the historical or physical probability measure, representing the actual likelihood of market events occurring.
\end{itemize}

Crucially, in a dynamic multi-period model, information arrives progressively over time. This flow of information is mathematically modeled by a \textbf{filtration} $\mathbb{F} = (\mathcal{F}_t)_{t=0, \dots, T}$. 
A filtration is an increasing sequence of sub-$\sigma$-algebras of $\mathcal{F}$, meaning that $\mathcal{F}_0 \subseteq \mathcal{F}_1 \subseteq \dots \subseteq \mathcal{F}_T = \mathcal{F}$. Intuitively, $\mathcal{F}_t$ represents the set of all information available to investors up to and including time $t$. Any decision made at time $t$ can only depend on information in $\mathcal{F}_t$; it cannot depend on the future.

\begin{remark}
At time $t=0$, we assume that the initial prices of all assets, $(S_0^0, S_0^1, \dots, S_0^d)$, are strictly observable and known. Because there is no uncertainty regarding the initial state of the market, the initial $\sigma$-algebra is trivial:
$$
\mathcal{F}_0 = \{\emptyset, \Omega\}
$$
This implies that any $\mathcal{F}_0$-measurable random variable is deterministically constant (almost surely).
\end{remark}

\subsection{Price Processes}

Let $S_t^i$ denote the price of the $i$-th asset at time $t$. The price evolution of asset $i$ over time is represented by a stochastic process $S^i = (S_t^i)_{t=0, \dots, T}$.

For the model to be economically sensible, the price $S_t^i$ must be known at time $t$. Mathematically, this means that the stochastic process $S^i$ is \textbf{adapted} to the filtration $\mathbb{F}$. That is, for every $t \in \{0, \dots, T\}$, the random variable $S_t^i$ is $\mathcal{F}_t$-measurable:
$$
\sigma\left(S_t^1, \dots, S_t^d\right) \subseteq \mathcal{F}_t, \quad \forall t \in \{0, \dots, T\}.
$$

\subsection{An Explicit Construction}

It is instructive to provide an explicit canonical construction for this abstract setup, focusing on a framework where the only source of randomness is the asset prices themselves.

\paragraph{The Riskless Asset:}
The riskless asset, denoted by $S^0$, models a bank account generating a constant, risk-free interest rate $r \geq 0$ per period. Assuming an initial price of $S_0^0 = 1$, the deterministic evolution of the riskless asset is defined as:
$$
S_t^0 := (1+r)^t, \quad t=0, \dots, T.
$$
This represents discrete compounding where 1 unit of currency invested at $t=0$ grows deterministically over time.

\paragraph{The Sample Space $\Omega$:}
Since we have $d$ risky assets traded over $T$ periods, a complete path of the strictly positive risky asset prices can be represented as a vector in $\mathbb{R}_+^{d \times T}$. We can explicitly define the sample space $\Omega$ as the Cartesian product:
$$
\Omega := \underbrace{\mathbb{R}_+^d \times \dots \times \mathbb{R}_+^d}_{T \text{ times}}.
$$
An element $\omega \in \Omega$ represents a specific price path: $\omega = (\omega_1, \dots, \omega_T)$, where each $\omega_t = (\omega_{t,1}, \dots, \omega_{t,d}) \in \mathbb{R}_+^d$ is the vector of prices for the $d$ assets at time $t$. The "universe" $\mathcal{F}$ is given by the Borel $\sigma$-algebra $\mathcal{B}(\Omega)$.

\paragraph{Asset Price Functions:}
We then define the price of asset $i$ at time $t$ as the projection mapping from the entire path space $\Omega$ to the specific value at time $t$:
$$
S_t^i : \Omega \to \mathbb{R}_+, \quad S_t^i(\omega) := \omega_{t,i}.
$$
Writing $S_t = (S_t^1, \dots, S_t^d)$, we have $S_t(\omega) = \omega_t$.

The natural choice for the filtration in this explicit formulation is the natural filtration generated by the price processes themselves. We define:
$$
\mathcal{F}_t := \sigma(S_1, \dots, S_t), \quad t = 1, \dots, T,
$$
with $\mathcal{F}_T = \mathcal{F}$. Under this setup, $\mathcal{F}_t$ embodies "the history of prices up to time $t$," and no exogenous information is considered.

\subsection{Discounted Processes}

A common simplification in pricing theory is to eliminate the effect of the deterministic interest rate by expressing all prices relative to the riskless asset $S^0$. This is equivalent to choosing the bank account as the \textbf{numéraire}.

We define the \textbf{discounted price processes} $\tilde{S}^i$ for $i = 0, \dots, d$:
$$
\tilde{S}_t^i := \frac{S_t^i}{S_t^0} = \frac{S_t^i}{(1+r)^t}, \quad t=0, \dots, T.
$$
In this discounted unit, the riskless asset's price is constant: $\tilde{S}_t^0 = 1$ for all $t$. By applying this transformation, we factor out the time value of money, dramatically simplifying the characterization of martingales and self-financing strategies.

\section{Self-Financing Strategies and Portfolios}

In a dynamic market, an investor continuously updates their portfolio by buying and selling assets according to a strategy. A trading strategy is formally defined by a $(d+1)$-dimensional stochastic process $\xi = (\xi^0, \xi^1, \dots, \xi^d)$, where:
\begin{itemize}
    \item $\xi_t^0$ represents the number of units held in the riskless asset between time $t-1$ and time $t$.
    \item $\xi_t^i$ (for $i=1, \dots, d$) represents the number of units held in the $i$-th risky asset between time $t-1$ and time $t$.
\end{itemize}

A critical requirement is that $\xi$ must be \textbf{predictable}. A process $\xi$ is predictable relative to the filtration $\mathbb{F}$ if for every $t \in \{1, \dots, T\}$, the random variable $\xi_t$ is $\mathcal{F}_{t-1}$-measurable. Economic intuition dictates this: the portfolio allocation $\xi_t$ chosen to be held during the interval $[t-1, t)$ must be decided at time $t-1$, using only the information available up to time $t-1$. It cannot foresee the prices that will be revealed at time $t$. For compactness, let us denote $\xi_t = (\xi_t^1, \dots, \xi_t^d)$. 

\begin{definition}[Self-Financing Strategy]
A predictable trading strategy $(\xi^0, \xi) = (\xi_t^0, \xi_t)_{t=1, \dots, T}$ taking values in $\mathbb{R}^{d+1}$ is called \textbf{self-financing} if it satisfies:
\begin{equation}
\label{eq:self_financing_condition}
\left(\xi_t^0, \xi_t\right) \cdot \left(S_t^0, S_t\right) = \left(\xi_{t+1}^0, \xi_{t+1}\right) \cdot \left(S_t^0, S_t\right) \quad \text{for all } t = 1, \dots, T-1.
\end{equation}
Equivalently:
$$
\xi_t^0 S_t^0 + \sum_{i=1}^d \xi_t^i S_t^i = \xi_{t+1}^0 S_t^0 + \sum_{i=1}^d \xi_{t+1}^i S_t^i
$$
\end{definition}

\textbf{Interpretation:} Equation \eqref{eq:self_financing_condition} perfectly captures the essence of a self-financing portfolio. 
At time $t$, before any rebalancing occurs, the investor holds the portfolio $\xi_t$ which was decided at $t-1$. The value of this portfolio evaluated at the newly revealed prices $S_t$ is the left-hand side: $(\xi_t^0, \xi_t) \cdot (S_t^0, S_t)$.
After observing the prices, the investor decides to rebalance their portfolio to the new allocation $\xi_{t+1}$, which will be held from time $t$ to $t+1$. The cost of assembling this new portfolio is exactly the right-hand side: $(\xi_{t+1}^0, \xi_{t+1}) \cdot (S_t^0, S_t)$.

The equality asserts that the value of the portfolio right before rebalancing exactly covers the cost of the newly assembled portfolio. There is no exogenous infusion of capital, nor is any capital withdrawn for consumption. All purchases are funded strictly by the sale of existing assets in the portfolio.

\begin{definition}[Portfolio Value Process]
The portfolio associated with a self-financing strategy $(\xi_0, \xi) = (\xi_t^0, \xi_t)_{t=1, \dots, T}$ has an adapted value process $(V_t^{\xi_0, \xi})_{t=0, \dots, T}$ described by:
\begin{align*}
V_0^{\xi_0, \xi} &= \left(\xi_1^0, \xi_1\right) \cdot \left(S_0^0, S_0\right) = \xi_1^0 + \sum_{i=1}^d \xi_1^i S_0^i, \\
V_t^{\xi_0, \xi} &= \left(\xi_t^0, \xi_t\right) \cdot \left(S_t^0, S_t\right) = \xi_t^0(1+r)^t + \sum_{i=1}^d \xi_t^i S_t^i, \quad t=1, \dots, T.
\end{align*}
\end{definition}

Because the strategy is self-financing, any fluctuation in the portfolio's value is purely a consequence of changes in asset prices, not changes in position volumes. We can formally calculate the variation of the portfolio's value from time $t$ to $t+1$:
\begin{align*}
V_{t+1}^{\xi_0, \xi} - V_t^{\xi_0, \xi} &= \left(\xi_{t+1}^0, \xi_{t+1}\right) \cdot \left(S_{t+1}^0, S_{t+1}\right) - \left(\xi_t^0, \xi_t\right) \cdot \left(S_t^0, S_t\right) \\
&= \left(\xi_{t+1}^0, \xi_{t+1}\right) \cdot \left(S_{t+1}^0, S_{t+1}\right) - \left(\xi_{t+1}^0, \xi_{t+1}\right) \cdot \left(S_t^0, S_t\right) \quad \text{(using the self-financing relation)} \\
&= \left(\xi_{t+1}^0, \xi_{t+1}\right) \cdot \left( (S_{t+1}^0, S_{t+1}) - (S_t^0, S_t) \right).
\end{align*}
This elegantly demonstrates that the change in portfolio value is completely attributed to the returns of the underlying assets $(\Delta S_{t+1})$ multiplied by the predictable allocations $(\xi_{t+1})$.

\subsection{Characterizations of Self-Financing Strategies}

A crucial mathematical property is how self-financing strategies interact with discounted prices. Working with discounted prices often significantly eases derivations.

\begin{proposition} \label{prop:self_financing_equiv}
Let $(\xi_t^0, \xi_t)_{t=1, \dots, T}$ be a trading strategy. The following claims are equivalent:
\begin{itemize}
    \item[(i)] The strategy $(\xi_t^0, \xi_t)_{t=1, \dots, T}$ is self-financing.
    \item[(ii)] The self-financing condition holds when evaluated with discounted prices:
    $$
    \left(\xi_t^0, \xi_t\right) \cdot \left(1, \tilde{S}_t\right) = \left(\xi_{t+1}^0, \xi_{t+1}\right) \cdot \left(1, \tilde{S}_t\right), \quad t=1, \dots, T-1.
    $$
    \item[(iii)] The discounted portfolio value $\tilde{V}_t^{\xi_0, \xi} := \frac{V_t^{\xi_0, \xi}}{S_t^0}$ can be expressed purely as stochastic discrete-time integral against the discounted risky asset price processes:
    $$
    \tilde{V}_t^{\xi_0, \xi} = V_0^{\xi_0, \xi} + \sum_{k=1}^t \xi_k \cdot (\tilde{S}_k - \tilde{S}_{k-1}), \quad t=0, \dots, T.
    $$
\end{itemize}
\end{proposition}

\begin{proof}
\textbf{(i) $\iff$ (ii):} Dividing the standard self-financing condition \eqref{eq:self_financing_condition} by the strictly positive observable scalar $S_t^0$ immediately yields the condition in (ii), and multiplying by $S_t^0$ reverses it.

\textbf{(i) $\iff$ (iii):} Assume the strict self-financing criteria holds. We compute the variation in the discounted portfolio value at time $t$:
\begin{align*}
\tilde{V}_t^{\xi_0, \xi} - \tilde{V}_{t-1}^{\xi_0, \xi} &= \left(\xi_t^0, \xi_t\right) \cdot \left(1, \tilde{S}_t\right) - \left(\xi_{t-1}^0, \xi_{t-1}\right) \cdot \left(1, \tilde{S}_{t-1}\right) \\
&= \left(\xi_t^0, \xi_t\right) \cdot \left(1, \tilde{S}_t\right) - \left(\xi_t^0, \xi_t\right) \cdot \left(1, \tilde{S}_{t-1}\right) \quad \text{(using property (ii) above)}\\
&= \left(\xi_t^0, \xi_t\right) \cdot \left(\left(1, \tilde{S}_t\right) - \left(1, \tilde{S}_{t-1}\right)\right) \\
&= \xi_t^0(1 - 1) + \xi_t \cdot (\tilde{S}_t - \tilde{S}_{t-1}) \\
&= \xi_t \cdot (\tilde{S}_t - \tilde{S}_{t-1}).
\end{align*}
Summing these differences consecutively from $k=1$ to $t$ transforms a telescoping sum into exactly the expression given in (iii):
$$
\tilde{V}_t^{\xi_0, \xi} = \tilde{V}_0^{\xi_0, \xi} + \sum_{k=1}^t \xi_k \cdot (\tilde{S}_k - \tilde{S}_{k-1})
$$
Noting that $\tilde{V}_0 = V_0 / S_0^0 = V_0 / 1 = V_0$, the result is proven. 

The reverse implication proceeds mechanically. Setting $\tilde{V}_t$ by the sum dictates exactly that $\tilde{V}_t - \tilde{V}_{t-1} = \xi_t \cdot (\tilde{S}_t - \tilde{S}_{t-1})$, which easily unwinds back into the definition of property (ii).
\end{proof}

The true significance of property (iii) is that the bank account position $(\xi_t^0)$ strictly vanishes from the representation of the discounted portfolio’s gains. This leads perfectly into the following corollary proposition:

\begin{proposition}
Let $(\xi_t)_{t=1, \dots, T}$ be an $\mathbb{F}$-predictable process taking values in $\mathbb{R}^d$ representing the holdings of risky assets. For every initial capital $x \in \mathbb{R}$, there exists a \textbf{unique} predictable process $(\xi_t^0)_{t=1, \dots, T}$ taking values in $\mathbb{R}$ such that the composite process $(\xi_t^0, \xi_t)_{t=1, \dots, T}$ forms a self-financing strategy starting with an initial portfolio value $V_0^{\xi_0, \xi} = x$.
\end{proposition}
Because the risky asset weights completely dictate the aggregate portfolio accumulation when combined with the self-financing mandate, the money-market position acts simply as a residual balancing term. The initial capital $x$ and the trades $\xi_t$ uniquely define the value of the portfolio at all times: 
$$
V_t = x \cdot S_t^0 + S_t^0 \sum_{k=1}^t \xi_k \cdot (\tilde{S}_k - \tilde{S}_{k-1}).
$$
To understand why $x$ is multiplied by $S_t^0$, recall from Proposition \ref{prop:self_financing_equiv} that the discounted portfolio value is bounded mathematically as $\tilde{V}_t = x + \sum_{k=1}^t \xi_k \cdot (\tilde{S}_k - \tilde{S}_{k-1})$. Multiplying both sides by the numéraire $S_t^0$ to obtain the exact actual value $V_t$ distributes the multiplication to $x$. Financially, the term $x \cdot S_t^0$ represents the initial capital capitalized to time $t$ by the risk-free rate (as if it had been entirely and passively invested in the bank account), because we assume $S_0^0 = 1$. The second term represents the accumulated trading gains from the risky assets, correctly capitalized to current time $t$. Consequently, we often safely denote a uniquely defined portfolio path strictly by its initial capital and risky allocations, writing $V^{x, \xi}$.
